\SourceLeaf{253}%
\noindent of proving the equality of volumes of symmetric solids, to which he replied that a similar proof could already be found in \textsc{Legendre}'s Geometry. I have yet to verify this.

\Needspace{7\baselineskip}
\begin{center}
[\textsc{Gerling} \emph{to} \textsc{Gauss}. \emph{Marburg, 21 July 1844.}]
\end{center}

\{\mbox{.\ .\ .\ .\ .\ .\ .} On account of Mr. \textsc{Eisenstein}'s remark about the volumes of symmetric solids, I consulted \textsc{Legendre}. I still have the edition of 1806 myself, but have now had an opportunity to examine \textsc{Crelle}'s translation (1844) of the 12th edition. Compared with 1806, I do indeed find some changes on these matters. --- \textsc{Legendre}, however, did not prove the proposition in question \textls{independently} of the method of exhaustion. He starts from the proposition that pyramids of equal base area and equal height are equal. --- He proves this by the method of exhaustion (in substance now just as in my \textsc{Lorenz} § 330, whereas earlier he followed Euclid more closely). He then divides the symmetric polyhedra into symmetric triangular pyramids and infers their equality by virtue of that proposition. The difference is therefore that \textsc{Legendre} does not carry the subdivision further to obtain 12 pyramids in each case, congruent in pairs. ---

I no longer remember how I learned the proposition in \textsc{Lorenz} § 330. Even if I discovered it independently, remarks by \textsc{Pfaff} certainly brought the discovery very close. I remember in particular that in (1809) he spoke to me of “Segner's axiom” and regretted that one could not penetrate further here.\}
